%%%%%%%%%%%%%%%%%%%%%%%%%%%%%%%%%%%%%%%%%%%%%%%%%%%%%%%%%%%%%%%%%%%%%%%%%%%%%%%%
% CHAPTER 3
%%%%%%%%%%%%%%%%%%%%%%%%%%%%%%%%%%%%%%%%%%%%%%%%%%%%%%%%%%%%%%%%%%%%%%%%%%%%%%%%

\chapter{Continuous-Time LTI Systems and Convolution}

\begin{tcolorbox}[colback=blue!5!white,colframe=blue!70!black,title=Learning Objectives]

After studying this chapter, the reader will be able to

\begin{itemize}
\item Define a continuous-time LTI system.
\item Understand impulse response and its significance.
\item Derive the convolution integral.
\item Compute the output of an LTI system using convolution.
\item Understand the graphical interpretation of convolution.
\item Use the properties of convolution in problem solving.
\end{itemize}

\end{tcolorbox}

%%%%%%%%%%%%%%%%%%%%%%%%%%%%%%%%%%%%%%%%%%%%%%%%%%%%%%%%%%%%%%%%%%%%%%%%%%%%%%%%

\section{Introduction}

A \textbf{Linear Time-Invariant (LTI) system} is the most important class of
systems in Signals and Systems because

\begin{itemize}
\item they are mathematically tractable,
\item they model many practical physical systems,
\item and their behavior is completely determined by a single signal called the
impulse response.
\end{itemize}

Examples of systems that can often be modeled as LTI systems include

\begin{itemize}
\item RC circuits,
\item RL circuits,
\item mechanical mass-spring systems,
\item communication channels,
\item and many filters used in signal processing.
\end{itemize}

%%%%%%%%%%%%%%%%%%%%%%%%%%%%%%%%%%%%%%%%%%%%%%%%%%%%%%%%%%%%%%%%%%%%%%%%%%%%%%%%

\section{Definition of an LTI System}

A system is LTI if it satisfies

\begin{enumerate}
\item \textbf{Linearity}
\item \textbf{Time invariance}
\end{enumerate}

%%%%%%%%%%%%%%%%%%%%%%%%%%%%%%%%%%%%%%%%%%%%%%%%%%%%%%%%%%%%%%%%%%%%%%%%%%%%%%%%

\subsection{Linearity}

For inputs \(x_1(t)\) and \(x_2(t)\),

\[
T\{a_1x_1(t)+a_2x_2(t)\}
=
a_1T\{x_1(t)\}+a_2T\{x_2(t)\}.
\]

%%%%%%%%%%%%%%%%%%%%%%%%%%%%%%%%%%%%%%%%%%%%%%%%%%%%%%%%%%%%%%%%%%%%%%%%%%%%%%%%

\subsection{Time Invariance}

If

\[
T\{x(t)\}=y(t),
\]

then

\[
T\{x(t-t_0)\}=y(t-t_0).
\]

%%%%%%%%%%%%%%%%%%%%%%%%%%%%%%%%%%%%%%%%%%%%%%%%%%%%%%%%%%%%%%%%%%%%%%%%%%%%%%%%

\section{Impulse Response}

%%%%%%%%%%%%%%%%%%%%%%%%%%%%%%%%%%%%%%%%%%%%%%%%%%%%%%%%%%%%%%%%%%%%%%%%%%%%%%%%

\subsection{Definition}

The response of an LTI system to a unit impulse input is called the
\textbf{impulse response}.

If

\[
x(t)=\delta(t),
\]

then

\[
\boxed{
y(t)=h(t)
}
\]

where \(h(t)\) is the impulse response.

%%%%%%%%%%%%%%%%%%%%%%%%%%%%%%%%%%%%%%%%%%%%%%%%%%%%%%%%%%%%%%%%%%%%%%%%%%%%%%%%

\subsection{Why is \(h(t)\) Important?}

For an LTI system, knowing \(h(t)\) is sufficient to determine the output for
\emph{any} input signal.

\[
\boxed{
\text{Input }x(t)
\quad \xrightarrow{\text{LTI}}
\quad
\text{Output }y(t)
}
\]

The relationship between \(x(t)\) and \(h(t)\) is called
\textbf{convolution}.

%%%%%%%%%%%%%%%%%%%%%%%%%%%%%%%%%%%%%%%%%%%%%%%%%%%%%%%%%%%%%%%%%%%%%%%%%%%%%%%%

\section{Impulse Decomposition of a Signal}

Any signal can be represented as a continuous sum of shifted impulses:

\[
\boxed{
x(t)=\int_{-\infty}^{\infty}x(\tau)\delta(t-\tau)\,d\tau
}
\]

This is obtained from the sifting property of the delta function.

%%%%%%%%%%%%%%%%%%%%%%%%%%%%%%%%%%%%%%%%%%%%%%%%%%%%%%%%%%%%%%%%%%%%%%%%%%%%%%%%

\section{Response to a Shifted Impulse}

For input

\[
\delta(t-\tau),
\]

the output is

\[
h(t-\tau).
\]

Because of linearity and time invariance,

\[
\boxed{
\delta(t-\tau)\ \longrightarrow\ h(t-\tau)
}
\]

%%%%%%%%%%%%%%%%%%%%%%%%%%%%%%%%%%%%%%%%%%%%%%%%%%%%%%%%%%%%%%%%%%%%%%%%%%%%%%%%

\section{Derivation of the Convolution Integral}

Starting from

\[
x(t)=\int_{-\infty}^{\infty}x(\tau)\delta(t-\tau)\,d\tau,
\]

apply the system operator:

\[
y(t)=T\{x(t)\}.
\]

Using linearity,

\[
y(t)=
\int_{-\infty}^{\infty}
x(\tau)\,T\{\delta(t-\tau)\}\,d\tau.
\]

Since

\[
T\{\delta(t-\tau)\}=h(t-\tau),
\]

we obtain

\[
\boxed{
y(t)=\int_{-\infty}^{\infty}x(\tau)h(t-\tau)\,d\tau
}
\]

This is the \textbf{convolution integral}.

%%%%%%%%%%%%%%%%%%%%%%%%%%%%%%%%%%%%%%%%%%%%%%%%%%%%%%%%%%%%%%%%%%%%%%%%%%%%%%%%

\section{Convolution Notation}

\[
\boxed{
y(t)=x(t)*h(t)
}
\]

where \(*\) denotes convolution.

Thus,

\[
\boxed{
x(t)*h(t)
=
\int_{-\infty}^{\infty}x(\tau)h(t-\tau)\,d\tau
}
\]

%%%%%%%%%%%%%%%%%%%%%%%%%%%%%%%%%%%%%%%%%%%%%%%%%%%%%%%%%%%%%%%%%%%%%%%%%%%%%%%%

\section{Graphical Interpretation}

Convolution can be understood through four steps:

%%%%%%%%%%%%%%%%%%%%%%%%%%%%%%%%%%%%%%%%%%%%%%%%%%%%%%%%%%%%%%%%%%%%%%%%%%%%%%%%

\subsection{Step 1: Flip}

Replace \(h(\tau)\) by

\[
h(-\tau).
\]

%%%%%%%%%%%%%%%%%%%%%%%%%%%%%%%%%%%%%%%%%%%%%%%%%%%%%%%%%%%%%%%%%%%%%%%%%%%%%%%%

\subsection{Step 2: Shift}

Shift by \(t\):

\[
h(t-\tau).
\]

%%%%%%%%%%%%%%%%%%%%%%%%%%%%%%%%%%%%%%%%%%%%%%%%%%%%%%%%%%%%%%%%%%%%%%%%%%%%%%%%

\subsection{Step 3: Multiply}

Multiply

\[
x(\tau)\,h(t-\tau).
\]

%%%%%%%%%%%%%%%%%%%%%%%%%%%%%%%%%%%%%%%%%%%%%%%%%%%%%%%%%%%%%%%%%%%%%%%%%%%%%%%%

\subsection{Step 4: Integrate}

Integrate over all \(\tau\):

\[
y(t)=\int x(\tau)h(t-\tau)\,d\tau.
\]

%%%%%%%%%%%%%%%%%%%%%%%%%%%%%%%%%%%%%%%%%%%%%%%%%%%%%%%%%%%%%%%%%%%%%%%%%%%%%%%%

\section{Convolution of Two Unit Steps}

Let

\[
x(t)=u(t),
\qquad
h(t)=u(t).
\]

Then,

\[
y(t)=\int_{-\infty}^{\infty}u(\tau)u(t-\tau)\,d\tau.
\]

%%%%%%%%%%%%%%%%%%%%%%%%%%%%%%%%%%%%%%%%%%%%%%%%%%%%%%%%%%%%%%%%%%%%%%%%%%%%%%%%

\subsection{Determine Limits}

The integrand is nonzero when

\[
\tau\ge0
\]

and

\[
t-\tau\ge0
\quad\Rightarrow\quad
\tau\le t.
\]

Therefore,

\[
0\le\tau\le t.
\]

%%%%%%%%%%%%%%%%%%%%%%%%%%%%%%%%%%%%%%%%%%%%%%%%%%%%%%%%%%%%%%%%%%%%%%%%%%%%%%%%

\subsection{Evaluate}

\[
y(t)=\int_0^t 1\,d\tau=t.
\]

Hence,

\[
\boxed{
u(t)*u(t)=t\,u(t)
}
\]

This result is extremely important.

%%%%%%%%%%%%%%%%%%%%%%%%%%%%%%%%%%%%%%%%%%%%%%%%%%%%%%%%%%%%%%%%%%%%%%%%%%%%%%%%

\section{Graphical Result}

\begin{center}
\begin{tikzpicture}[scale=0.9]
\draw[->] (-0.5,0) -- (5,0) node[right] {$t$};
\draw[->] (0,-0.5) -- (0,4) node[above] {$y(t)$};

\draw[thick] (-0.5,0) -- (0,0);
\draw[thick] (0,0) -- (4,3.2);
\node[right] at (3,2.6) {$t\,u(t)$};
\end{tikzpicture}
\end{center}

%%%%%%%%%%%%%%%%%%%%%%%%%%%%%%%%%%%%%%%%%%%%%%%%%%%%%%%%%%%%%%%%%%%%%%%%%%%%%%%%

\section{Convolution with an Impulse}

Let

\[
x(t)*\delta(t).
\]

Using the convolution integral,

\[
y(t)=
\int_{-\infty}^{\infty}
x(\tau)\delta(t-\tau)\,d\tau.
\]

By the sifting property,

\[
\boxed{
x(t)*\delta(t)=x(t)
}
\]

Thus, the impulse acts as the identity element for convolution.

%%%%%%%%%%%%%%%%%%%%%%%%%%%%%%%%%%%%%%%%%%%%%%%%%%%%%%%%%%%%%%%%%%%%%%%%%%%%%%%%

\section{Important Convolution Pairs}

\begin{table}[H]
\centering
\caption{Common convolution results}
\begin{tabular}{|p{5cm}|p{5cm}|}
\hline
\textbf{Signals} & \textbf{Result} \\
\hline
\(x(t)*\delta(t)\) & \(x(t)\) \\
\hline
\(u(t)*u(t)\) & \(t\,u(t)\) \\
\hline
\(u(t)*t\,u(t)\) & \(\dfrac{t^2}{2}u(t)\) \\
\hline
\(e^{-at}u(t)*u(t)\) & \(\dfrac{1-e^{-at}}{a}u(t)\) \\
\hline
\end{tabular}
\end{table}

%%%%%%%%%%%%%%%%%%%%%%%%%%%%%%%%%%%%%%%%%%%%%%%%%%%%%%%%%%%%%%%%%%%%%%%%%%%%%%%%

\section{Properties of Convolution}

%%%%%%%%%%%%%%%%%%%%%%%%%%%%%%%%%%%%%%%%%%%%%%%%%%%%%%%%%%%%%%%%%%%%%%%%%%%%%%%%

\subsection{Commutative}

\[
\boxed{
x(t)*h(t)=h(t)*x(t)
}
\]

%%%%%%%%%%%%%%%%%%%%%%%%%%%%%%%%%%%%%%%%%%%%%%%%%%%%%%%%%%%%%%%%%%%%%%%%%%%%%%%%

\subsection{Associative}

\[
\boxed{
[x_1(t)*x_2(t)]*x_3(t)
=
x_1(t)*[x_2(t)*x_3(t)]
}
\]

%%%%%%%%%%%%%%%%%%%%%%%%%%%%%%%%%%%%%%%%%%%%%%%%%%%%%%%%%%%%%%%%%%%%%%%%%%%%%%%%

\subsection{Distributive}

\[
\boxed{
x(t)*[h_1(t)+h_2(t)]
=
x(t)*h_1(t)+x(t)*h_2(t)
}
\]

These properties greatly simplify convolution calculations.

%%%%%%%%%%%%%%%%%%%%%%%%%%%%%%%%%%%%%%%%%%%%%%%%%%%%%%%%%%%%%%%%%%%%%%%%%%%%%%%%

\section{Worked Example 3.1}

Find

\[
y(t)=u(t)*u(t).
\]

\textbf{Solution}

\[
y(t)=\int_{-\infty}^{\infty}u(\tau)u(t-\tau)\,d\tau
\]

The valid region is

\[
0\le\tau\le t.
\]

Therefore,

\[
y(t)=\int_0^t 1\,d\tau=t.
\]

Hence,

\[
\boxed{
y(t)=t\,u(t)
}
\]

%%%%%%%%%%%%%%%%%%%%%%%%%%%%%%%%%%%%%%%%%%%%%%%%%%%%%%%%%%%%%%%%%%%%%%%%%%%%%%%%

\begin{tcolorbox}[title=Quick Recall,colback=yellow!10]

\[
y(t)=x(t)*h(t)
\]

\[
y(t)=\int_{-\infty}^{\infty}x(\tau)h(t-\tau)\,d\tau
\]

\[
u(t)*u(t)=t\,u(t)
\]

\[
x(t)*\delta(t)=x(t)
\]

Flip \(\rightarrow\) Shift \(\rightarrow\) Multiply \(\rightarrow\) Integrate

\end{tcolorbox}
%%%%%%%%%%%%%%%%%%%%%%%%%%%%%%%%%%%%%%%%%%%%%%%%%%%%%%%%%%%%%%%%%%%%%%%%%%%%%%%%
%%%%%%%%%%%%%%%%%%%%%%%%%%%%%%%%%%%%%%%%%%%%%%%%%%%%%%%%%%%%%%%%%%%%%%%%%%%%%%%%
\section{Graphical Convolution in Detail}

Consider

\[
x(t)=u(t),
\qquad
h(t)=u(t)-u(t-1).
\]

The impulse response is a rectangular pulse of width 1.

%%%%%%%%%%%%%%%%%%%%%%%%%%%%%%%%%%%%%%%%%%%%%%%%%%%%%%%%%%%%%%%%%%%%%%%%%%%%%%%%

\subsection{Step 1: Flip}

\[
h(-\tau)=u(-\tau)-u(-\tau-1).
\]

%%%%%%%%%%%%%%%%%%%%%%%%%%%%%%%%%%%%%%%%%%%%%%%%%%%%%%%%%%%%%%%%%%%%%%%%%%%%%%%%

\subsection{Step 2: Shift}

\[
h(t-\tau)=u(t-\tau)-u(t-\tau-1).
\]

%%%%%%%%%%%%%%%%%%%%%%%%%%%%%%%%%%%%%%%%%%%%%%%%%%%%%%%%%%%%%%%%%%%%%%%%%%%%%%%%

\subsection{Step 3: Overlap Region}

The overlap exists when

\[
\tau\ge0
\]

and

\[
t-1\le\tau\le t.
\]

Hence,

\[
\max(0,t-1)\le\tau\le t.
\]

%%%%%%%%%%%%%%%%%%%%%%%%%%%%%%%%%%%%%%%%%%%%%%%%%%%%%%%%%%%%%%%%%%%%%%%%%%%%%%%%

\subsection{Step 4: Integrate}

For \(0<t<1\),

\[
y(t)=\int_0^t 1\,d\tau=t.
\]

For \(t\ge1\),

\[
y(t)=\int_{t-1}^{t}1\,d\tau=1.
\]

Therefore,

\[
\boxed{
y(t)=
\begin{cases}
0, & t<0\\
t, & 0\le t<1\\
1, & t\ge1
\end{cases}
}
\]

%%%%%%%%%%%%%%%%%%%%%%%%%%%%%%%%%%%%%%%%%%%%%%%%%%%%%%%%%%%%%%%%%%%%%%%%%%%%%%%%

\section{Convolution of Rectangular Pulses}

Let

\[
x(t)=u(t)-u(t-T),
\]

\[
h(t)=u(t)-u(t-T).
\]

Both are pulses of width \(T\).

The result is a triangular waveform:

\[
\boxed{
y(t)=
\begin{cases}
0, & t<0\\
t, & 0\le t<T\\
2T-t, & T\le t<2T\\
0, & t\ge2T
\end{cases}
}
\]

%%%%%%%%%%%%%%%%%%%%%%%%%%%%%%%%%%%%%%%%%%%%%%%%%%%%%%%%%%%%%%%%%%%%%%%%%%%%%%%%

\subsection{Graph}

\begin{center}
\begin{tikzpicture}[scale=0.9]
\draw[->] (-0.5,0) -- (5.5,0) node[right] {$t$};
\draw[->] (0,-0.5) -- (0,3) node[above] {$y(t)$};

\draw[thick] (0,0) -- (2,2) -- (4,0);

\node[below] at (2,0) {$T$};
\node[below] at (4,0) {$2T$};
\end{tikzpicture}
\end{center}

This result is frequently asked in GATE.

%%%%%%%%%%%%%%%%%%%%%%%%%%%%%%%%%%%%%%%%%%%%%%%%%%%%%%%%%%%%%%%%%%%%%%%%%%%%%%%%

\section{Convolution with an Exponential}

Let

\[
x(t)=u(t),
\qquad
h(t)=e^{-at}u(t),
\qquad a>0.
\]

%%%%%%%%%%%%%%%%%%%%%%%%%%%%%%%%%%%%%%%%%%%%%%%%%%%%%%%%%%%%%%%%%%%%%%%%%%%%%%%%

\subsection{Convolution Integral}

\[
y(t)=\int_{-\infty}^{\infty}
u(\tau)e^{-a(t-\tau)}u(t-\tau)\,d\tau.
\]

The valid region is

\[
0\le\tau\le t.
\]

Hence,

\[
y(t)=\int_0^t e^{-a(t-\tau)}\,d\tau.
\]

%%%%%%%%%%%%%%%%%%%%%%%%%%%%%%%%%%%%%%%%%%%%%%%%%%%%%%%%%%%%%%%%%%%%%%%%%%%%%%%%

\subsection{Evaluation}

\[
y(t)=e^{-at}\int_0^t e^{a\tau}\,d\tau
\]

\[
=e^{-at}\left[\frac{e^{a\tau}}{a}\right]_0^t
\]

\[
=\frac{1-e^{-at}}{a}.
\]

Therefore,

\[
\boxed{
u(t)*e^{-at}u(t)
=
\frac{1-e^{-at}}{a}u(t)
}
\]

%%%%%%%%%%%%%%%%%%%%%%%%%%%%%%%%%%%%%%%%%%%%%%%%%%%%%%%%%%%%%%%%%%%%%%%%%%%%%%%%

\section{Output of an LTI System}

For any input,

\[
\boxed{
y(t)=x(t)*h(t)
}
\]

Thus, the entire behavior of an LTI system is contained in \(h(t)\).

%%%%%%%%%%%%%%%%%%%%%%%%%%%%%%%%%%%%%%%%%%%%%%%%%%%%%%%%%%%%%%%%%%%%%%%%%%%%%%%%

\section{Causality of CT-LTI Systems}

A continuous-time LTI system is causal if and only if

\[
\boxed{
h(t)=0,\qquad t<0
}
\]

%%%%%%%%%%%%%%%%%%%%%%%%%%%%%%%%%%%%%%%%%%%%%%%%%%%%%%%%%%%%%%%%%%%%%%%%%%%%%%%%

\subsection{Example}

\[
h(t)=e^{-t}u(t)
\]

Since \(u(t)=0\) for \(t<0\),

\[
\boxed{
\text{The system is causal}
}
\]

%%%%%%%%%%%%%%%%%%%%%%%%%%%%%%%%%%%%%%%%%%%%%%%%%%%%%%%%%%%%%%%%%%%%%%%%%%%%%%%%

\section{Stability of CT-LTI Systems}

A CT-LTI system is BIBO stable if

\[
\boxed{
\int_{-\infty}^{\infty}|h(t)|\,dt<\infty
}
\]

%%%%%%%%%%%%%%%%%%%%%%%%%%%%%%%%%%%%%%%%%%%%%%%%%%%%%%%%%%%%%%%%%%%%%%%%%%%%%%%%

\subsection{Example}

\[
h(t)=e^{-t}u(t)
\]

\[
\int_0^\infty e^{-t}dt=1<\infty
\]

Hence,

\[
\boxed{
\text{The system is stable}
}
\]

%%%%%%%%%%%%%%%%%%%%%%%%%%%%%%%%%%%%%%%%%%%%%%%%%%%%%%%%%%%%%%%%%%%%%%%%%%%%%%%%

\section{Impulse Response and Differential Equations}

Many physical systems are described by differential equations.

%%%%%%%%%%%%%%%%%%%%%%%%%%%%%%%%%%%%%%%%%%%%%%%%%%%%%%%%%%%%%%%%%%%%%%%%%%%%%%%%

\subsection{First-Order System}

\[
\boxed{
\frac{dy(t)}{dt}+ay(t)=x(t)
}
\]

To find \(h(t)\), apply

\[
x(t)=\delta(t).
\]

Then,

\[
\frac{dh(t)}{dt}+ah(t)=\delta(t).
\]

The solution is

\[
\boxed{
h(t)=e^{-at}u(t)
}
\]

This is the impulse response of a first-order stable system.

%%%%%%%%%%%%%%%%%%%%%%%%%%%%%%%%%%%%%%%%%%%%%%%%%%%%%%%%%%%%%%%%%%%%%%%%%%%%%%%%

\section{RC Circuit Example}

Consider a series RC circuit with output across the capacitor.

The governing equation is

\[
RC\frac{dy(t)}{dt}+y(t)=x(t).
\]

Comparing with the standard form,

\[
a=\frac{1}{RC}.
\]

Therefore,

\[
\boxed{
h(t)=\frac{1}{RC}e^{-t/(RC)}u(t)
}
\]

%%%%%%%%%%%%%%%%%%%%%%%%%%%%%%%%%%%%%%%%%%%%%%%%%%%%%%%%%%%%%%%%%%%%%%%%%%%%%%%%

\section{Step Response}

The step response is the output for

\[
x(t)=u(t).
\]

Using convolution,

\[
s(t)=u(t)*h(t).
\]

For the RC circuit,

\[
s(t)=
\int_0^t
\frac{1}{RC}e^{-(t-\tau)/(RC)}\,d\tau.
\]

Evaluating,

\[
\boxed{
s(t)=\left(1-e^{-t/(RC)}\right)u(t)
}
\]

This is the familiar charging equation of a capacitor.

%%%%%%%%%%%%%%%%%%%%%%%%%%%%%%%%%%%%%%%%%%%%%%%%%%%%%%%%%%%%%%%%%%%%%%%%%%%%%%%%

\section{Convolution Summary}

\begin{table}[H]
\centering
\caption{Important convolution results}
\begin{tabular}{|p{6cm}|p{5cm}|}
\hline
\textbf{Convolution} & \textbf{Result} \\
\hline
\(x(t)*\delta(t)\) & \(x(t)\) \\
\hline
\(u(t)*u(t)\) & \(t\,u(t)\) \\
\hline
\(u(t)*t\,u(t)\) & \(\dfrac{t^2}{2}u(t)\) \\
\hline
\(u(t)*e^{-at}u(t)\) & \(\dfrac{1-e^{-at}}{a}u(t)\) \\
\hline
Rectangular pulse * rectangular pulse & Triangular pulse \\
\hline
\end{tabular}
\end{table}

%%%%%%%%%%%%%%%%%%%%%%%%%%%%%%%%%%%%%%%%%%%%%%%%%%%%%%%%%%%%%%%%%%%%%%%%%%%%%%%%

\section{Problem-Solving Strategy}

%%%%%%%%%%%%%%%%%%%%%%%%%%%%%%%%%%%%%%%%%%%%%%%%%%%%%%%%%%%%%%%%%%%%%%%%%%%%%%%%

\subsection{Step 1}

Identify \(x(t)\) and \(h(t)\).

%%%%%%%%%%%%%%%%%%%%%%%%%%%%%%%%%%%%%%%%%%%%%%%%%%%%%%%%%%%%%%%%%%%%%%%%%%%%%%%%

\subsection{Step 2}

Write

\[
y(t)=\int x(\tau)h(t-\tau)\,d\tau.
\]

%%%%%%%%%%%%%%%%%%%%%%%%%%%%%%%%%%%%%%%%%%%%%%%%%%%%%%%%%%%%%%%%%%%%%%%%%%%%%%%%

\subsection{Step 3}

Determine the overlap region.

%%%%%%%%%%%%%%%%%%%%%%%%%%%%%%%%%%%%%%%%%%%%%%%%%%%%%%%%%%%%%%%%%%%%%%%%%%%%%%%%

\subsection{Step 4}

Set the correct limits of integration.

%%%%%%%%%%%%%%%%%%%%%%%%%%%%%%%%%%%%%%%%%%%%%%%%%%%%%%%%%%%%%%%%%%%%%%%%%%%%%%%%

\subsection{Step 5}

Evaluate the integral.

%%%%%%%%%%%%%%%%%%%%%%%%%%%%%%%%%%%%%%%%%%%%%%%%%%%%%%%%%%%%%%%%%%%%%%%%%%%%%%%%

\subsection{Step 6}

Express the answer using unit-step functions if required.

%%%%%%%%%%%%%%%%%%%%%%%%%%%%%%%%%%%%%%%%%%%%%%%%%%%%%%%%%%%%%%%%%%%%%%%%%%%%%%%%

\section{Common GATE Mistakes}

\begin{enumerate}
\item Forgetting to flip \(h(\tau)\) before shifting.
\item Using incorrect overlap limits.
\item Forgetting the unit-step multiplier.
\item Confusing convolution with ordinary multiplication.
\item Assuming convolution is not commutative.
\item Ignoring the causality condition \(h(t)=0\) for \(t<0\).
\item Forgetting the absolute value in the stability integral.
\end{enumerate}

%%%%%%%%%%%%%%%%%%%%%%%%%%%%%%%%%%%%%%%%%%%%%%%%%%%%%%%%%%%%%%%%%%%%%%%%%%%%%%%%

\section{Chapter Summary}

In this chapter, we studied continuous-time LTI systems and showed that the
output of an LTI system is obtained by convolving the input with the impulse
response.

We derived the convolution integral from the impulse decomposition of a signal
and learned the graphical interpretation of convolution through the
flip-shift-multiply-integrate procedure.

We also studied

\begin{itemize}
\item convolution of step signals,
\item convolution of rectangular pulses,
\item convolution with exponential signals,
\item causality and stability conditions,
\item and impulse responses obtained from differential equations.
\end{itemize}

These concepts form the foundation for Fourier analysis and frequency-domain
representation of systems.

%%%%%%%%%%%%%%%%%%%%%%%%%%%%%%%%%%%%%%%%%%%%%%%%%%%%%%%%%%%%%%%%%%%%%%%%%%%%%%%%

\begin{tcolorbox}[title=One-Minute Revision,colback=blue!5]

\[
y(t)=x(t)*h(t)
\]

\[
y(t)=\int_{-\infty}^{\infty}x(\tau)h(t-\tau)\,d\tau
\]

\[
u(t)*u(t)=t\,u(t)
\]

\[
u(t)*e^{-at}u(t)=\frac{1-e^{-at}}{a}u(t)
\]

Causal:

\[
h(t)=0,\quad t<0
\]

Stable:

\[
\int_{-\infty}^{\infty}|h(t)|dt<\infty
\]

RC impulse response:

\[
h(t)=\frac{1}{RC}e^{-t/(RC)}u(t)
\]

RC step response:

\[
s(t)=\left(1-e^{-t/(RC)}\right)u(t)
\]

\end{tcolorbox}
%%%%%%%%%%%%%%%%%%%%%%%%%%%%%%%%%%%%%%%%%%%%%%%%%%%%%%%%%%%%%%%%%%%%%%%%%%%%%%%%